% SEGMENT OLP-0064-S01
% Part: first-order-logic
% Chapter: proof-systems
% Section: introduction

\documentclass[../../../include/open-logic-section]{subfiles}

\begin{document}

\iftag{FOL}
      {\olfileid{fol}{prf}{int}}
      {\olfileid{pl}{prf}{int}}

\olsection{அறிமுகம்}

தருக்கவியல்களில் பொதுவாக ஒரு பொருண்மையியலும் !!a{derivation}
முறைமையும் உள்ளன. பொருண்மையியல் மெய்மை, நிறைவுறத்தக்கதன்மை,
செல்லுபடித்தன்மை, பின்விளைவு போன்ற கருத்துகளைப் பற்றியது.
!!{derivation} முறைமைகளின் நோக்கம், பின்விளைவையும் செல்லுபடித்தன்மையையும்
நிறுவுவதற்கு முற்றிலும் தொடரமைப்பு சார்ந்த ஒரு முறையை வழங்குவதாகும்.
இத்தகைய முறைமையில் !!a{derivation} என்பது முடிவுறு தொடரமைப்புப் பொருள்;
பொதுவாக அது !!{sentence}s அல்லது !!{formula}s இன் தொடர் (அல்லது வேறொரு
முடிவுறு ஒழுங்கமைப்பு) ஆகும். எந்தவொரு !!{sentence}s அல்லது
!!{formula}s தொடரோ ஒழுங்கமைப்போ ``சரியானதா'' என்பதை இயந்திர முறையில்
சரிபார்க்க முடியும் என்ற பண்பு நல்ல !!{derivation} முறைமைகளுக்கு உண்டு.

% SEGMENT OLP-0064-S02
முதல்-வரிசைத் தருக்கவியலுக்கான மிக எளிய (வரலாற்றிலும் முதலில் தோன்றிய)
!!{derivation} முறைமைகள் \emph{அடிகோள் சார்ந்தவை}. அத்தகைய முறைமையில்
!!{formula}s இன் ஒரு தொடர் !!a{derivation} ஆகக் கருதப்பட வேண்டுமெனில்,
அதிலுள்ள ஒவ்வொரு !!{formula}வும் நிலையான ``அடிகோள்கள்'' கணத்தில்
ஒன்றாகவோ, தொடரில் அதற்கு முன் வரும் !!{formula}s இலிருந்து நிலையான
எண்ணிக்கையிலான ``அனுமான விதிகள்'' ஒன்றால் பின்வருவதாகவோ இருக்க வேண்டும்.
மேலும், !!a{formula} ஓர் அடிகோளா என்பதையும், ஓர் அனுமான விதியால் பிற
!!{formula}s இலிருந்து அது சரியாகப் பின்வருகிறதா என்பதையும் இயந்திர
முறையில் சரிபார்க்க முடியும். அடிகோள் சார்ந்த !!{derivation} முறைமைகளை
விவரிப்பதும் அவற்றை மேற்கோட்பாட்டளவில் கையாள்வதும் எளிது; ஆனால் அவற்றின்
!!{derivation}s ஐ வாசித்துப் புரிந்துகொள்வதும் உருவாக்குவதும் கடினம்.

% SEGMENT OLP-0064-S03
முழுமையான !!{derivation}s ஐ எளிதாக அமைக்கவும், அமைத்த பிறகு
!!{derivation}s ஐ எளிதாகப் புரிந்துகொள்ளவும் வேறு !!{derivation} முறைமைகள்
உருவாக்கப்பட்டன.
இயல்பான வருவித்தல், அட்டவணைமர நிறுவல்கள் என்றும் அழைக்கப்படும் மெய்மை
மரங்கள், தொடரணிக் கணியம் ஆகியவை எடுத்துக்காட்டுகள். சில !!{derivation}
முறைமைகள் குறிப்பாக இயந்திரமயமாக்கலை நோக்கமாகக் கொண்டு வடிவமைக்கப்பட்டவை;
எடுத்துக்காட்டாக, தீர்வாக்க முறையை மென்பொருளில் செயல்படுத்துவது எளிது
(ஆனால் அதன் !!{derivation}s ஐப் புரிந்துகொள்வது கிட்டத்தட்ட இயலாது).
இந்த மற்ற !!{derivation} முறைமைகளில் பெரும்பாலானவை !!{derivation}s ஐ
!!{formula}s இன் தொடர்களாக அல்லாமல் மரங்களாகக் காட்டுகின்றன. இதனால்
!!a{derivation}வின் எந்தப் பகுதிகள் எந்த மற்றப் பகுதிகளைச் சார்ந்துள்ளன
என்பதைக் காண்பது எளிதாகிறது.

% SEGMENT OLP-0064-S04
ஆகவே முதல்-வரிசைத் தருக்கவியல் போன்ற ஒரு தருக்கவியலுக்கு, வெவ்வேறு
!!{derivation} முறைமைகள் !!a{sentence} ஒரு \emph{தேற்றம்} என்பது
என்னவென்றும், !!a{sentence} சில மற்ற வாக்கியங்களிலிருந்து
!!{derivable} என்பது என்னவென்றும் வெவ்வேறு துல்லிய விளக்கங்களைத் தரும்.
அதை எவ்வாறு செய்தாலும் (அடிகோள் சார்ந்த !!{derivation}s, இயல்பான
வருவித்தல்கள், தொடரணி !!{derivation}s, மெய்மை மரங்கள், தீர்வாக்க
மறுப்புகள் ஆகியவற்றின் வழியாக), அந்த உறவுகள் செல்லுபடித்தன்மை மற்றும்
பின்விளைவு என்னும் பொருண்மைக் கருத்துகளுடன் பொருந்த வேண்டும். ``$!A$~ஒரு
தேற்றம்'' என்பதற்கு $\Proves !A$ என்றும், ``$\Gamma$ இலிருந்து $!A$
!!{derivable}'' என்பதற்கு ``$\Gamma \Proves !A$'' என்றும் எழுதுவோம்.
$\Proves$ எவ்வாறு வரையறுக்கப்பட்டாலும், அது $\Entails$ உடன் பின்வருமாறு
பொருந்த வேண்டும்:
\begin{enumerate}
\item $\Proves !A$ என்பது, அப்போதும் அப்போது மட்டுமே, $\Entails !A$
\item $\Gamma \Proves !A$ என்பது, அப்போதும் அப்போது மட்டுமே,
  $\Gamma \Entails !A$
\end{enumerate}

% SEGMENT OLP-0064-S05
மேலுள்ளவற்றின் ``அப்போது மட்டுமே'' திசை \emph{சரித்தன்மை} எனப்படுகிறது.
!!{derivability} பின்விளைவை (அல்லது செல்லுபடித்தன்மையை) உறுதிசெய்தால்,
!!^a{derivation} முறைமை சரித்தன்மையுடையது. பயன்படத்தக்க ஒவ்வொரு
!!{derivation} முறைமையும் சரித்தன்மையுடையதாக இருக்க வேண்டும்;
சரித்தன்மையற்ற !!{derivation} முறைமைகள் முற்றிலும் பயனற்றவை. ஏனெனில்
செல்லுபடித்தன்மை அல்லது பின்விளைவுக்குத் தொடரமைப்பு சார்ந்த உத்தரவாதம்
வழங்குவதே !!a{derivation}வின் முழு நோக்கம். நாம் வழங்கும் !!{derivation}
முறைமைகளின் சரித்தன்மையை நிறுவுவோம்.

% SEGMENT OLP-0064-S06
மறுதிசையான ``எனில்'' திசையும் முக்கியமானது; அது \emph{நிறைவு}
எனப்படுகிறது. $!A$~செல்லுபடியாகும் போதெல்லாம் $!A$~ஒரு தேற்றம் என்பதை
காட்டவும், $\Gamma \Entails !A$ எனும் போதெல்லாம் $\Gamma \Proves !A$
என்பதைக் காட்டவும் நிறைவான !!{derivation} முறைமை போதிய வலிமை கொண்டது.
நிறைவை நிறுவுவது கடினமானது; சில தருக்கவியல்களுக்கு நிறைவான
!!{derivation} முறைமைகளே இல்லை. முதல்-வரிசைத் தருக்கவியலுக்கு அத்தகைய
முறைமை உண்டு. கூர்ட் கேடல் தமது 1929 ஆய்வேட்டில் முதல்-வரிசைத்
தருக்கவியலுக்கான !!a{derivation} முறைமையின் நிறைவை முதன்முதலில்
நிறுவினார்.

% SEGMENT OLP-0064-S07
!!{derivation} முறைமைகளோடு தொடர்புடைய மற்றொரு கருத்து
\emph{முரணின்மை}. !!{sentence}s இன் ஒரு கணத்திலிருந்து எதை
வேண்டுமானாலும் !!{derive} இயலுமானால் அது முரணுடையது என்றும், இல்லையெனில்
முரணற்றது என்றும் கூறப்படும். முரணுடைமை என்பது நிறைவுறத்தகாமைக்கு
தொடரமைப்பு சார்ந்த இணைக்கருத்து. நிறைவுறத்தகாத கணங்களைப் போலவே,
!!{sentence}s இன் முரணுடைய கணங்களும் நல்ல கோட்பாடுகளாக அமையாது; அவற்றில்
அடிப்படையான குறைபாடு உள்ளது. !!{sentence}s இன் முரணற்ற கணங்கள் மெய்யாகவோ
பயனுள்ளதாகவோ இருக்க வேண்டியதில்லை; ஆனால் தருக்கப் பயன்பாட்டிற்கான அந்தக்
குறைந்தபட்சத் தகுதியையாவது அவை அடைகின்றன. வெவ்வேறு !!{derivation}
முறைமைகளில் !!{sentence}s கணங்களின் முரணின்மைக்கான குறிப்பிட்ட வரையறை
மாறலாம். ஆனால் $\Proves$ போலவே, முரணின்மையும் அதன் பொருண்மை இணையான
நிறைவுறத்தக்கதன்மையும் ஒன்றுபட வேண்டும். $\Gamma$ முரணற்றது என்பது,
அப்போதும் அப்போது மட்டுமே, அது நிறைவுறத்தக்கது என்பதாக எப்போதும் இருக்க
வேண்டும். இங்கு ``அப்போது மட்டுமே'' திசை நிறைவைக் குறிக்கும் (முரணின்மை
நிறைவுறத்தக்கதன்மையை உறுதிசெய்கிறது); ``எனில்'' திசை சரித்தன்மையைக்
குறிக்கும் (நிறைவுறத்தக்கதன்மை முரணின்மையை உறுதிசெய்கிறது). உண்மையில்,
செவ்வியல் முதல்-வரிசைத் தருக்கவியலுக்கு சரித்தன்மை மற்றும் நிறைவின்
இரண்டு வடிவங்களும் ஒன்றுக்கொன்று நிகரானவை.

\end{document}
